\documentclass[margin,line]{res}
\usepackage[T1]{fontenc}
\usepackage{lmodern}
\usepackage{microtype}
\usepackage{hyperref}
\usepackage[brief]{cvbib}
\addbibresource{publications.bib}

\oddsidemargin -.5in
\evensidemargin -.5in
\textwidth=6.0in
\setlength{\emergencystretch}{3em}
\renewcommand*{\sectionfont}{\normalfont\scshape} % res.cls default is bold

% Role line: title left, date right. If both don't fit on one line, the date
% moves to the next line, still flush right. No page break right after it.
\newcommand*{\cvrole}[2]{\emph{#1}\unskip\nobreak\hfill\penalty50\hskip1em\hbox{}\nobreak\hfill\mbox{\textbf{#2}}}

% No page break inside a list, or between it and the line introducing it,
% so a role and its bullets stay on one page.
\makeatletter
\newenvironment{list2}{
  \begin{list}{$\bullet$}{%
      \@beginparpenalty=\@M \@itempenalty=\@M
      \setlength{\itemsep}{0in}
      \setlength{\parsep}{0in} \setlength{\parskip}{0in}
      \setlength{\topsep}{0in} \setlength{\partopsep}{0in} 
      \setlength{\leftmargin}{0.2in}}}{\end{list}}
\makeatother
      
\hypersetup{%
  pdftitle={Eric L. Denovellis, CV},
  pdfauthor={Eric L. Denovellis},
  colorlinks=true,
  urlcolor=blue,
  linkcolor=blue
}


\begin{document}

\name{Eric L. Denovellis, Ph.D. \vspace*{.1in}}

\begin{resume}
\section{Contact Information}
\vspace{.05in}
\begin{tabular}{@{}p{1.6in}p{4in}}
San Francisco, CA & \textit{Email:}  \href{mailto:eric.denovellis@ucsf.edu}{eric.denovellis@ucsf.edu}\\            
& \textit{Website:}  \href{https://www.edenovellis.com}{edenovellis.com}\\
& \textit{GitHub:}  \href{https://github.com/edeno}{github.com/edeno}\\
& \textit{Google Scholar:}  \href{https://scholar.google.com/citations?user=vuENbkoAAAAJ}{Eric Denovellis}
\end{tabular}

\section{Summary}
Research scientist with 10+ years developing real-time neural decoding algorithms and scalable data infrastructure for systems neuroscience. Expert in state-space models, point processes, and Bayesian inference for neural signal processing. Builder of open-source tools used by 15+ research labs.

\section{Skills} 
\begin{list2}
\item \textbf{Programming:} Python, JAX, Matlab, R, JavaScript, SQL
\item \textbf{ML/Statistics:} State-space models, point processes, Bayesian inference, neural decoding, time series analysis, spectral analysis, Kalman filtering
\item \textbf{Real-time Systems:} Online clusterless decoding, closed-loop neural feedback, GPU-accelerated signal processing
\item \textbf{Data Infrastructure:} DataJoint, Spyglass, NWB (Neurodata Without Borders), pipeline development, database design
\item \textbf{Visualization:} Interactive web-based tools, D3.js, matplotlib, scientific figure design
\end{list2}

\section{Research Experience}
\textbf{University of California, San Francisco}

\cvrole{Research Scientist}{September 2019 – Present}\\*
Frank Lab, Department of Physiology
\begin{list2}
\item Developed marked point process state-space models for real-time decoding of hippocampal replay, enabling closed-loop neurofeedback experiments (Neuron, 2025)
\item Extended these models to decode replay across multiple spatial environments (COSYNE, 2023)
\item Co-led development of Spyglass, now used by 15+ neuroscience labs (eLife, 2026)
\item Added GPU acceleration to spectral\_connectivity for terabyte-scale data (JOSS, 2022)
\item Led team of 2 computational scientists; mentored 5+ trainees on statistical methods and software development
\item Collaborated with hardware engineers on a real-time clusterless decoder with 6 ms temporal resolution and $<$50 ms median computational latency on 32+ tetrodes (eNeuro, 2025)
\end{list2}

\textbf{Boston University}

\cvrole{Postdoctoral Researcher}{September 2016 – August 2019}\\*
Mathematics and Statistics Department, Eden and Kramer Labs
\begin{list2}
\item Designed novel switching state-space models for characterizing hippocampal replay dynamics
\item Developed hybrid point process filters for real-time trajectory decoding from population spiking activity
\item Created spectral\_connectivity, an open-source package for spectral and connectivity analysis
\item Created interactive web-based visualization tools for multi-dimensional electrophysiology data exploration
\end{list2}

\cvrole{Graduate Researcher}{September 2009 – September 2016}\\*
Bullock Lab (BU) and Miller Lab (MIT)
\begin{list2}
\item Discovered synchronous oscillatory mechanisms in prefrontal cortex that coordinate context switching (Neuron, 2012; 250+ citations)
\item Applied generalized additive models and point process methods to characterize neural encoding of cognitive variables
\end{list2}

\section{Software}

\href{https://github.com/LorenFrankLab/spyglass}{\textbf{Spyglass}} (co-first author) — Data analysis framework for reproducible neuroscience research. Integrates NWB, DataJoint, SpikeInterface, and DeepLabCut. \textbf{121 stars, 62 forks.}

\href{https://github.com/Eden-Kramer-Lab/spectral_connectivity}{\textbf{spectral\_connectivity}} — GPU-accelerated multitaper spectral estimation and connectivity analysis. \textbf{137 stars, 49 forks.}

\href{https://github.com/LorenFrankLab/non_local_detector}{\textbf{non\_local\_detector}} — GPU-accelerated Bayesian state-space models in JAX for decoding hippocampal replay from sorted or clusterless spikes.

\href{https://github.com/Eden-Kramer-Lab/replay_trajectory_classification}{\textbf{replay\_trajectory\_classification}} — State-space models for decoding hippocampal replay trajectories and classifying their dynamics from sorted or clusterless spikes. \textbf{53 stars, 20 forks.}

\href{https://github.com/Eden-Kramer-Lab/ripple_detection}{\textbf{ripple\_detection}} — Sharp wave ripple and high synchrony event detection from LFP and multiunit activity. \textbf{47 stars, 16 forks.}

\href{https://github.com/NeurophysVis/SpectraVis}{\textbf{SpectraVis}} — Interactive D3.js application for exploring brain connectivity. \textbf{37 stars, 12 forks.}

\section{Selected Publications}
Full publication list: \href{https://scholar.google.com/citations?user=vuENbkoAAAAJ}{Google Scholar} \hfill \emph{* Co-first authors}
\vspace{.1in}

\nocite{lee2026spyglass,chu2025realtimedecoder,coulter2025closedloop,joshi2023stepping,denovellis2022spectral,denovellis2021replay,buschman2012synchronous}
\printbibliography[heading=none,env=cvbib]

\section{Talks \&\\ Teaching}
\begin{list2}
\item \textbf{Invited talks:} Spring Hippocampal Symposium, Verona, NeuroAI workshop (upcoming, 2027); Neuroscience Frontiers, Columbia University (upcoming, 2026); NCBS, India (2026); UCL Gatsby/SWC (2024); University of Washington (2024); UBC (2024)
\item \textbf{Teaching:} Visiting Instructor, NCBS Bangalore (2026), four-week computational neuroscience course; Guest Lecturer, UCSF (2024)
\end{list2}

\section{Education}
\textbf{Boston University} — Ph.D., Computational Neuroscience, 2016\\
\textbf{University of California, Santa Barbara} — B.S., Mathematics; B.A., Philosophy, 2007

\end{resume}
\end{document}
